\documentclass[11pt]{article}
\usepackage{mymacros}
\usepackage{mytheorems}
\usepackage{textcomp,geometry,graphicx,enumerate,fullpage}
\usepackage{booktabs}
\usepackage{setspace}
\usepackage{hyperref}
\usepackage{mathrsfs} % for mathscr
\usepackage{cleveref}
\usepackage{todonotes}
\DeclareMathOperator{\tr}{tr}

\newcommand{\suppfun}[1]{h_{#1}}
\newcommand{\error}{e}
\newcommand{\Error}{\mathbf{E}_A}
\newcommand{\APsd}{\mathbf{Q}_A}

\title{Hidden Convexity}

\date{\today}

\newcommand{\cA}{\mathcal{A}}

\DeclareMathOperator{\vecf}{vec}

\begin{document}
\maketitle

\section{Outline}

In this paper, we introduce a semi-automated way to analyse the global landscape
of nonconvex optimization problems.  More precisely, the goal is to investigate
the existence of local minima that are not global, which we called
\emph{spurious}.  The approach requires finding a \emph{convex} problem for dual
solutions provide certificate of global optimality of specific solution.

This convex formulation can be used to either find spurious local minima or aid
the development of a proof of global optimality for all local minima.  We
illustrate the methodology on two nonconvex problems: the Burer-Monteiro
formulation for Semidefinite Programming.  We show how the approach can
streamline the analysis.

For Burer-Monteiro, we start with the general case where we show how the method
can be used to find existing counterexamples found by
\cite{bhojanapalli2018smoothed}.  For the particular case of Burer-Monteiro
applied to the Sum-of-Squares problem, we show how the approach find the proof
of \cite{legat2023lowrank} and the counterexamples of
\cite{blekherman2024Spurious}.

We argue that such hidden convexity approach should provide a way for the
optimization community to leverage its expertise in convex modeling to the
analysis of the nonconvex formulation that are prevalent nowadays.  Moreover,
once these hidden convex problems are found, it significantly streamline the
proof development, by providing guidance through the search of these complicated
global landscape.

A similar approach is in \cite[(3.2)]{zhang2024improved}.  The key difference is
that their $\phi$ represents the whole problem so they optimize of the class of
all SDPs while we just optimize over the right-hand side of the constant of the
SDP.

\section{Background}
Let $S^n$ be the space of $n\times n$ symmetric matrices, and $S^n_r$ be the
space of rank-$r$ symmetric PSD matrices.  Let $\cA: \R^n \rightarrow \R^m$ be a
quadratic map defined by $A_1, \ldots, A_m \in S^n$:
\begin{align*}
  \cA(x) = (x^\top A_1 x, \ldots, x^\top A_m x).
\end{align*}
More generally, we can define
$A: S^n \rightarrow \R^m$:
\begin{align*}
  A(X) = (\dotp{A_1, X}, \ldots, \dotp{A_m, X}).
\end{align*}
We say that $A$ (resp. $\cA$) satisfies \emph{hidden convexity} if
$\braces{A(x) \mid X \in S_1^n} = \braces{\cA(x) \mid x \in \R^n}$ is a convex
set. More generally, we can also consider the set
$\braces{A(x) \mid X \in S_r^n}$ for $r < n$. Below we list a few examples of
hidden convexity.

\paragraph{Dines' Theorem} When $m=2$, Dines \cite{dines1941mapping} proved that
$\braces{(x^\top A_1 x, x^\top A_2 x) \mid x \in \R^n}$, which is
$ \braces{(\dotp{A_1, X}, \dotp{A_2, X}) \mid X \in S_1^n}$, is convex.

\paragraph{Barvinok-Pataki Bound} When $r \ge \floor{(\sqrt{8m+1}-1)/2}$, the
set $\braces{A(x) \mid X \in S_r^n}$ is convex. This is known as the
Barvinok-Pataki
\cite{BarvinokFeasibilitytestingsystems1993,PatakiRankExtremeMatrices1998}
bound. Dines' theorem arises as a special case of this bound, when $m = 2$.

\paragraph{Sum of Squares Polynomials} Let $b(x) \in \R^n$ be a basis of
$\Rx[p, d]$, the set of polynomials in $p$ variables of degree at most
$d$. Suppose $A(X)$ is the map extracting the coefficients of the polynomial
$b(x)^\top X b(x)$. Then if $r$ is at least the Pythagoras number of
$\Rx[p, d]$, the set $\braces{A(x) \mid X \in S_r^n}$ is equal to the sum of
squares cone $\Sx[p, 2d]$, which is convex.

\section{Landscape of Burer-Monteiro Factorization}
Consider the following problem: Given $A: S^n \rightarrow \R^m$ and
$b \in \R^m$, find $X \in S^n_r$ so that $A(X) = b$. We consider the following
quadratically-penalized \footnote{The squared norm can be replaced by any convex
  positive definite function and the same analysis should hold.} objective:
\begin{align} \label{eq:obj}
  \min_{U \in \R^{n \times r}} f_b(U) := \norm{A(UU^\top) - b}^2.
\end{align}
If $U$ is a second-order critical point (SOCP) of \eqref{eq:obj}, the following
conditions must hold:
\begin{align}
  \label{eq:grad}
  \dotp{A(UV^\top), A(UU^\top) - b} \sim \nabla f_b(U)(V) &= 0,\quad \text{for all } V \in \R^{n \times r}. \\
  \label{eq:hess}
  \dotp{A(VV^\top), A(UU^\top) - b} + 2\norm{A(UV^\top)}^2 = \nabla^2 f_b(U)(V, V) &\ge 0,\quad \text{for all } V \in \R^{n \times r}.
\end{align}
In order to prove this for all SOCPs $U$, we want to show that the two sets
\begin{align*}
  G_A(U) &:= \braces{c \in \R^m \mid \nabla f_c(U)(V) = 0,\, \forall\, V \in \R^{n \times r}}, \\
  H_A(U) &:= \braces{c \in \R^m \mid \nabla^2 f_c(U)(V, V) \ge 0,\, \forall\, V \in \R^{n \times r}},
\end{align*}
only intersect at the point $b$. As we will see, this depends on the interplay
between the quadratic map $Q(U) = A(UU^\top)$ and the linear map
$A_U(V) = A(UV^\top)$. It is more convenient to work with affine shifts of
$G_A(U)$ and $H_A(U)$, shifting $Q(U)$ to the origin:
\begin{align*}
  \textbf{G}_A(U)
  &:= \braces{\error \in \R^m \mid \dotp{A_U(V), \error} = 0,\, \forall\, V \in \R^{n \times r}}, \\
  \textbf{H}_A(U)
  &:= \braces{\error \in \R^m \mid \dotp{Q(V), \error} + 2 \norm{A_U(V)}^2 \ge 0,\, \forall\, V \in \R^{n \times r}}.
\end{align*}
We have a nice characterization of $\textbf{G}_A(U)$ in terms of the linear map
$A_U$: It is the orthogonal complement of $\imag(A_U)$. In other words,
\begin{align*}
  \textbf{G}_A(U) = \kernel(A_U^*).
\end{align*}
We wish to show that there are no spurious SOCPs for \eqref{eq:obj}, for a
particular fixed $U \in \R^{n\times r}$. For the simplest case, first consider
the case when $Q(U)$ is in the interior of $\APsd$.
\begin{proposition}
  If $Q(U)$ is in the interior of $\APsd$ and
  \begin{align*}
    H_A(U) \cap G_A(U) \ne \braces{Q(U)},
  \end{align*}
  then there exists $b \in \APsd$ where $U$ is a spurious SOCP for $f_b$
  (i.e. $U$ is SOCP but not global minimum of $f_b$).
  \begin{proof}
    Since $Q(U)$ is in the interior of $\APsd$ and $H_A(U) \cap G_A(U)$ is not a
    singleton, we can find $b \in H_A(U) \cap G_A(U)$ (ensuring that $U$ is a
    SOCP of $f_b$) where $b \in \APsd$ and $b \ne Q(U)$ (ensuring that $U$ is
    not a global minimum of $f_b$).
  \end{proof}
\end{proposition}

If $Q(U)$ is on the boundary of $\APsd$, the main difficulty arises from
characterizing $b$ for which $U$ is \emph{not} a global minimum of $f_b$. In
particular, there are $b$ for which $f_b(U) > 0$ but $U$ is still a global
minimum. This occurs when $Q(U)$ is the \emph{projection} of $b$ onto the convex
set $\APsd$. We can account for these cases with the strengthened proposition:
\begin{proposition}
  Given any $U \in \R^{n \times r}$, if there exists $c \in \APsd$ and
  $b \in H_A(U) \cap G_A(U)$ so that
  \begin{align*}
    \dotp{Q(U) - c, Q(U) - b} > 0,
  \end{align*}
  then $U$ is a spurious SOCP of $f_b$.
  \begin{proof}
    Since $b \in H_A(U) \cap G_A(U)$, $U$ is a SOCP of $f_b$. Using the
    variational characterization of projection onto the convex set $\APsd$: $U$
    is the minimum of $f_b(U) = \norm{Q(U) - b}^2$ if for all $c \in \APsd$,
    $\dotp{Q(U) - c, Q(U) - b} \le 0$. Thus $\dotp{Q(U) - c, Q(U) - b} > 0 $ for
    some $c \in \APsd$ if and only if $U$ not a global minimum of $f_b$.
  \end{proof}
\end{proposition}

In summary, to show that there are no spurious SOCPs for \eqref{eq:obj}, we need
to show that for all $U \in \R^{n \times r}$ satisfying \eqref{eq:grad} and
\eqref{eq:hess}, $H_A(U) \cap G_A(U) = \braces{b}$.  For example, if
$\Sigma = \braces{Q(U) \mid U \in \R^{n \times r}}$ is a pointed convex cone,
when $U = 0$ it is clear that these two sets only intersect at the origin.

\section{Counterexample SDP and Duality}
We can determine whether a fixed $U \in \R^{n \times r}$ is a spurious SOCP for
a particular $f_b(U)$ by conducting a search over all $b \in \R^m$. Let
\begin{align*}
  \APsd &:= \{ Q(V)\mid V \in \R^{n \times r} \},\\
  \Error(U) &:= \{ Q(U) - Q(V) \mid V \in \R^{n \times r} \} = Q(U) - \APsd.
\end{align*}

\begin{proposition}
  \label{prop:onepoint}
  $\textbf{G}_A(U) \cap \textbf{H}_A(U) \cap \Error(U) \ne \{0\}$ if and
  only if $U$ is a spurious SOCP for $f_b(U)$ for some $b \in \APsd$.

  \begin{proof}
    If $a \in \textbf{G}_A(U)\cap \textbf{H}_A(U) \cap \Error(U)$ is
    non-zero, let $b = Q(U) - a$. Then $U$ is a SOCP of $f_b$ by definition and
    $f_b(U) = \norm{a}^2 > 0$, but since $a = Q(U) - Q(V)$ for some $V$,
    $b = Q(V)$ and $f_b(V) = 0$ is the global minimum.

    On the other hand, if $U$ is a spurious SOCP of $f_b$ where $b = Q(V)$ for
    some $V$, then $Q(U) - b \in \textbf{G}_A(U)\cap \textbf{H}_A(U)$ but
    $f_b(U) > 0$. This implies that $Q(U) - b \ne 0$.
  \end{proof}
\end{proposition}

When $\APsd$ is a convex semidefinite-representable set, i.e. there is
hidden convexity, this search can be performed by solving a SDP. In particular,
hidden convexity means that $\APsd$ is the image of the PSD cone under
the linear map $A$:
\begin{align*}
  \APsd = \{ A(X) \mid X \in S_+^n \}.
\end{align*}

Let $H^1_A: \R^m \rightarrow \R^{nr \times nr}$ be the linear map defined by
$\vecf(V)^T H^1_A(a) \vecf(V) := \dotp{A(VV^\top), a}$ and
$H^2_{A, U} \in \R^{nr \times nr}$ be the quadratic form defined by
$\vecf(V)^T H^2_{A, U} \vecf(V) := 2\norm{A_U(V)}^2$.

We now reformulate
\[ \textbf{G}_A(U) \cap \textbf{H}_A(U) \cap \Error(U) \ne \{0\} \]
into
\[ \exists c \in \R^m \text{ s.t. } \suppfun{\textbf{G}_A(U) \cap \textbf{H}_A(U) \cap \Error(U)}(c) > 0. \]
where $\suppfun{S}$ is the support function of a set $S$.
Note that there if the set is nonempty then the set of $c$ such that the support function is nonzero is
at least a halfspace.
So picking a random $c$ already gives a 50\% chance of detecting a counterexample.

This support function evaluation be formulated as follows:
\begin{align*}
  \max_{e} \, & \dotp{-c, \error} &
  \min_{\lambda, \nu, Y} \, & \dotp{\nu, Q(U)} + \dotp{Y, H^2_{A, U}} \\
  \text{s.t. } & A_U^*(\error) = 0  &
  \text{s.t. } & A_U(\lambda) + {H^1_A}^*(Y) - c = \nu \\
         & H^1_A(\error) + H^2_{A, U} \succeq 0 &
         & Y \succeq 0\\
         & Q(U) - \error \in \APsd &
         & \nu \in \APsd^*
\end{align*}
or
\begin{align}
  \label{eq:progb}
  \max_{b} \, & \dotp{-c, Q(U) - b} &
  \min_{\lambda, P} \, & \dotp{c + A_U(\lambda) + {H^1_A}^*(P), Q(U)} + \dotp{P, H^2_{A, U}} \\
  \notag
  \text{s.t. } & A_U^*(Q(U) - b) = 0  &
  \text{s.t. } & A_U(\lambda) + {H^1_A}^*(P) + c \in \APsd^*\\
  \notag
         & H^1_A(Q(U) - b) + H^2_{A, U} \succeq 0 &
         & P \succeq 0\\
  \notag
         & b \in \APsd\\
\end{align}
If we expand $\APsd$ (noting that $\APsd^* = (A(\Psd))^* = A^{-*}(\Psd)$ where $A^{-*}$ is the preimage of the adjoint (TODO cite rockafellar)), we get
\begin{align*}
  \max_{X} \, & \dotp{-c, Q(U) - A(X)} &
  \min_{\lambda, P} \, & \dotp{c + A_U(\lambda) + {H^1_A}^*(P), Q(U)} + \dotp{P, H^2_{A, U}} \\
  \text{s.t. } & A_U^*(Q(U) - A(X)) = 0  &
  \text{s.t. } & A^*(A_U(\lambda)) + A^*({H^1_A}^*(P)) + A^*(c) \succeq 0\\
         & H^1_A(Q(U) - A(X)) + H^2_{A, U} \succeq 0 &
         & P \succeq 0 \\
         & X \succeq 0
\end{align*}
Note that, as $\dotp{Q(U), Q(U) - b} = 0$, we have $\dotp{-c, Q(U) - b} = \dotp{-c', Q(U) - b}$ where $c' = c - Q(U) \dotp{c, Q(U)} / \dotp{Q(U), Q(U)}$.
As $\dotp{c', Q(U)} = 0$, we further have $\dotp{-c', Q(U) - b} = \dotp{c', b}$.
The equivalence with the following program is due to the following proposition.
\begin{proposition}
  Consider a convex cone $K$ and a point $x \in K$.
  If $x$ is in the interior of $K$ or if $x$ has a unique dual point $x^*$ then
  for all $y \in \R^n$ such that $\dotp{x^*, y} \le 0$,
  there exists $\lambda \in \R_+$ such that
  $y + \lambda x \in K$.
  \begin{proof}
    We can simply take $\lambda = \inf \{\dotp{z, y} / \dotp{z, x} \mid z \in K^*, \dotp{z, x} > 0\}$.
  \end{proof}
\end{proposition}
Given a feasible solution of \cref{eq:progb},
we have $\dotp{Q(U) - b, Q(U) - b} > 0$ hence $c = b$ is a feasible solution of the following program.

The projection formulation, given $c \in \APsd$,
\begin{align*}
  \max_b \, & \dotp{Q(U) - c, Q(U) - b} &
  \min_{P, \lambda} \, & \dotp{P, H^2_{A, U}} \\
  \text{s.t. } & A_U^*(Q(U) - b) = 0 &
  \text{s.t. } & c - Q(U) = A_U(\lambda) + {H^1_A}^*(P) \\
         & H^1_A(Q(U) - b) + H^2_{A, U} \succeq 0 &
         & P \succeq 0
\end{align*}


Then we can find a
counterexample if there exists $c \in \R^m$ so that the following SDP has a
non-zero objective value:
\begin{align*}
  \max_e \, & \dotp{c, e} \\
  \text{s.t. } & A_U^*(e) = 0 \\
         & H^1_A(e) + H^2_{A, U} \succeq 0 \\
         &  A^*(Q(U) - e) \succeq 0
\end{align*}
We can homogenize by using $\sqrt{\gamma}U$ which gives
$H^2_{A, \sqrt{\gamma}U} = \gamma H^2_{A, U}$,
$Q(\sqrt{\gamma}U) = \gamma Q(U)$,
\begin{align*}
  \max_b \, & \dotp{c, b} \\
  \text{s.t. } & A_U^*(b) = 0 \\
         & H^1_A(b) + \gamma H^2_{A, U} \succeq 0 \\
         &  A^*(\gamma Q(U) - b) \succeq 0
\end{align*}
Now it's homogeneous in $(\gamma, b)$ so we can do wlog
\begin{align*}
  \text{s.t. } & \dotp{c, b} = 1 \\
  & A_U^*(b) = 0 \\
         & H^1_A(b) + \gamma H^2_{A, U} \succeq 0 \\
         &  A^*(\gamma Q(U) - b) \succeq 0
\end{align*}
Taking the dual of this SDP we get:
\begin{align*}
  \min_{X, Y, \lambda} \, & \dotp{A(X), Q(U)} + \dotp{Y, H^2_{A, U}} \\
  \text{s.t. } & A(X) = c + A_U(\lambda) + {H^1_A}^*(Y)  \\
                          & X, Y \succeq 0
\end{align*}

Another way is to certify that $U$ is not a spurious SOCP for any $b \in \R^m$
is to show that for all $c \in \APsd$, the following SDP has objective
value 0 (note that the objective is at least 0 because $b = Q(U)$ is a feasible
solution):
\begin{align*}
  (P_1) = \max_b \, & \dotp{Q(U) - c, Q(U) - b} \\
  \text{s.t. } & A_U^*(Q(U) - b) = 0 \\
         & H^1_A(Q(U) - b) + H^2_{A, U} \succeq 0
\end{align*}
This is because of the variational characterization of projection onto the
convex set $\APsd$: $U$ is the minimum of $f_b(U) = \norm{Q(U) - b}^2$ if for
all $c \in \APsd$, $\dotp{Q(U) - c, Q(U) - b} \le 0$. Thus the objective is
positive for some $c \in \APsd$ (i.e.  $\dotp{Q(U) - c, Q(U) - b} > 0 $) if and
only if $U$ is a SOCP but not a global minimum of $f_b$.

Note that, with a limit argument, we can replace $c \in \mathcal{Q}_A$ with $c$
being in the tangent space at $U$.

Notice that, because $A_U^*(Q(U) - b) = 0$, we have $\dotp{Q(U), Q(U) - b} = 0$.
So the objective can be rewritten as $\dotp{-c, Q(U) - b}$.  The existence of a
$c \in \APsd$ such that $\dotp{-c, Q(U) - b}$ is positive is equivalent to the
existence of $c \in \APsd$ such that $\dotp{c, Q(U) - b}$ is negative.  So there
should be a separating hyperplane $h$ separating $Q(U) - b$ to the cone
$K = \{f | \dotp{c, f} \ge 0, \forall c \in \APsd \}$.  This cone is equivalent
to $K = \{f | \dotp{Q, A^*(f)} \ge 0, \forall Q \text{ is PSD}\}$.  So
$A^*(K) = \Psd$ is the PSD cone.  So the dual of $K$ is $K^* = A^{-1}(\Psd)$
where $A^{-1}$ is the preimage of $A$.  So $h \in K^*$ is equivalent to
$A(h) \in \Psd$.

As $\APsd$ is the image of the PSD cone through $A$, the nonpositivity of that
scalar product for all $c$ is equivalent to $A^*(Q(U) - b) \succeq 0$.  So
$(P_1) \le 0$ is equivalent to the existence of PSD $X$ such that
$\langle A(X), e \rangle > 0$ where $\error = Q(U) - b$.

Making the substitution $\error = Q(U) - b$ and taking the dual, we
get
\begin{align*}
  \min_{P, \lambda} \, & \dotp{P, H^2_{A, U}} \\
  \text{s.t. } & c - Q(U) = A_U(\lambda) + {H^1_A}^*(P) \\
                      & P \succeq 0
\end{align*}
If strong duality holds, the primal and dual have the same objective value of 0,
implying that $\dotp{P, H^2_{A, U}} = 0$ and hence $P = VV^\top$ where
$V \in \ker(A_U)$.  This corresponds to the certificate form of our proof for
univariate sum of squares using the simpler condition (C2). However, as we have
seen, strong duality may not always hold as the cones involved may not be
closed. Therefore in the full proof we need a limiting argument to deal with the
closure issues.

\bibliographystyle{amsalpha}
\bibliography{MyLibraryTEX}

\end{document}
